\documentclass[11pt,a4paper]{article}
\usepackage[T1]{fontenc}
\usepackage[utf8]{inputenc}
\usepackage[french]{babel}
\usepackage[margin=2.4cm]{geometry}
\usepackage{amsmath,amssymb}
\usepackage{booktabs}
\usepackage{microtype}
\usepackage[hidelinks]{hyperref}

\newcommand{\R}{\mathbb{R}}
\newcommand{\norm}[1]{\lVert #1 \rVert}
\newcommand{\inner}[2]{\langle #1, #2\rangle}
\newcommand{\T}{^{\top}}

\title{\textbf{CIFAR-10 — Exercices 6 \& 7}\\[2pt]
\large Rapport et justification des choix}
\author{Projet FTML}
\date{\today}

\begin{document}
\maketitle
\vspace{-1.2em}

\begin{abstract}
\noindent On classe les images CIFAR-10 sans réseau de neurones, avec des
méthodes classiques d'optimisation et d'algèbre linéaire (SVM, noyaux, PCA,
$K$-means). L'exercice~6 traite la classification supervisée (SVM, puis passage à
l'échelle par Nyström) ; l'exercice~7 apprend des représentations non supervisées,
évaluées par sonde linéaire. Toutes les valeurs numériques proviennent du notebook
\texttt{cifar\_ex6\_7.ipynb}.
\end{abstract}

\section{Jeu de données et problème}
\paragraph{Présentation (CIFAR-10).} Le jeu de données est CIFAR-10, téléchargé
depuis \url{https://www.cs.toronto.edu/~kriz/cifar.html} (A.~Krizhevsky, 2009). Il
contient $60\,000$ petites photographies couleur de $32\times32$ pixels, réparties
en $10$ classes d'objets du quotidien : avion, automobile, oiseau, chat, cerf,
chien, grenouille, cheval, bateau, camion. Chaque image est fournie comme un
vecteur de $d=32\times32\times3=3072$ entiers dans $[0,255]$ (les trois canaux RVB
mis bout à bout, dans cet ordre), stocké dans des fichiers \texttt{pickle}. Les
classes sont \textbf{équilibrées} ($6\,000$ images chacune) et le découpage
officiel donne $50\,000$ images d'entraînement et $10\,000$ de test. Point moins
évident pour qui ne connaît pas ces données : aucune variable n'a de sens prise
isolément (un pixel seul n'apprend rien), l'information est dans les
\emph{corrélations spatiales} entre pixels voisins ; de plus certaines classes sont
visuellement proches (chat/chien, automobile/camion), ce qui annonce des confusions.

\paragraph{Problème.} On résout une \textbf{classification multi-classes} :
prédire l'étiquette $y\in\{0,\dots,9\}$ (la classe de l'objet) à partir des seules
$3072$ intensités de pixels $x$, sans information annexe. C'est une brique de base
de la vision par ordinateur (tri et étiquetage automatique d'images, modération de
contenu, indexation visuelle) ; CIFAR-10 sert surtout de banc d'essai standard pour
\emph{comparer} des méthodes, ce qui est précisément notre objectif.

\section{Cadre et protocole}
On note les données $\{(x_i,y_i)\}_{i=1}^n$, $x_i\in\R^{d}$, $d=3072$. Le SVM à
noyau \emph{exact} demande une matrice de Gram $n\times n$ ($\sim$20\,Go pour
$n=50\,000$) : on l'évalue donc sur un sous-échantillon $10\,000/2\,000$, tandis que
les méthodes linéaires, Nyström, PCA et $K$-means exploitent les données complètes.

\paragraph{Risque réel vs risque empirique.} On cherche un classifieur $f$ de petit
\emph{risque réel} $R(f)=\mathbb E_{(x,y)\sim p}\big[\ell(f(x),y)\big]$, espérance de
la perte sur la vraie distribution $p$ (inconnue), ici avec la perte $0/1$. Comme
$p$ est inaccessible, on ne dispose que du \emph{risque empirique}
$R_n(f)=\tfrac1n\sum_{i=1}^n \ell(f(x_i),y_i)$, moyenne sur l'échantillon, que les
modèles minimisent (éventuellement régularisé). L'écart $R(f)-R_n(f)$ est l'erreur
de \emph{généralisation} : un modèle qui colle au train ($R_n$ petit) mais
généralise mal ($R$ grand) \emph{sur-apprend}. Le jeu de test, jamais utilisé pour
l'apprentissage ni le réglage, fournit un estimateur non biaisé de $R(f)$ : c'est
pourquoi on ne le consulte qu'une fois, et l'accuracy test rapportée estime $1-R(f)$.

\paragraph{Préparation des données.} Les pixels sont \textbf{standardisés}
(centrage/réduction par variable) avant les modèles à distance ($k$-NN, SVM) : sans
cela, les coordonnées de plus forte variance domineraient le calcul des distances et
du noyau. Pour la PCA on centre simplement les données (la covariance n'a de sens
que centrée).

\paragraph{Hyperparamètres.} Pour le SVM RBF, on sélectionne le couple
$(C,\gamma)$ par \textbf{validation croisée} (\texttt{GridSearchCV}, $3$ plis
stratifiés) sur la grille $C\in\{1,10,100\}$, $\gamma\in\{\text{scale},10^{-3},
10^{-2}\}$. Une recherche exhaustive sur $50\,000$ images avec noyau exact étant
trop coûteuse, cette CV se fait sur un sous-échantillon d'entraînement
($3\,000$ images) ; le meilleur couple est ensuite ré-entraîné sur les $10\,000$
images, puis évalué \emph{une seule fois} sur le test. Pour les méthodes passant à
l'échelle (Nyström, $K$-means), trop lourdes pour une grille complète, on conserve
des valeurs usuelles. Dans tous les cas, le jeu de test ne sert jamais à choisir un
modèle ni un hyperparamètre.

\section{Exercice 6 — Classification supervisée}

\subsection{Baselines}
On situe trois modèles linéaires de référence (sur pixels bruts) :
\begin{itemize}\itemsep2pt
  \item \textbf{$k$-NN} : vote majoritaire des $k$ plus proches voisins ; aucune
  optimisation, illustration du compromis biais–variance.
  \item \textbf{Régression softmax} : $p(y=c\mid x)=e^{w_c\T x}/\sum_j e^{w_j\T x}$,
  minimisation de l'entropie croisée $-\sum_i\log p(y_i\mid x_i)$, fonction
  \emph{convexe} à gradient explicite.
  \item \textbf{SVM linéaire} : perte hinge régularisée
  $\tfrac12\norm{w}^2+C\sum_i\max(0,1-y_i(w\T x_i+b))$ (cf.~\ref{sec:svm}).
\end{itemize}

\subsection{SVM à marge maximale : dérivation}\label{sec:svm}
\paragraph{Primal (marge dure).} Pour des données séparables, $y_i\in\{\pm1\}$,
l'hyperplan $w\T x+b=0$ a une marge $2/\norm{w}$ ; la maximiser revient à
\begin{equation}
\min_{w,b}\ \tfrac12\norm{w}^2
\quad\text{s.c.}\quad y_i(w\T x_i+b)\ge 1 \ \ \forall i .
\end{equation}
\paragraph{Marge souple.} Pour des données non séparables, on relâche par des
écarts $\xi_i\ge0$ :
\begin{equation}
\min_{w,b,\xi}\ \tfrac12\norm{w}^2 + C\sum_i\xi_i
\quad\text{s.c.}\quad y_i(w\T x_i+b)\ge 1-\xi_i,\ \xi_i\ge0 ,
\end{equation}
dont l'optimum donne $\xi_i=\max(0,1-y_i(w\T x_i+b))$ : la perte \emph{hinge}.
\paragraph{Lagrangien.} Avec $\alpha_i\ge0$ (marge) et $\mu_i\ge0$ ($\xi_i\ge0$),
\[
\mathcal L=\tfrac12\norm{w}^2+C\sum_i\xi_i
-\sum_i\alpha_i\big[y_i(w\T x_i+b)-1+\xi_i\big]-\sum_i\mu_i\xi_i .
\]
L'annulation des dérivées donne les conditions stationnaires
\begin{equation}
w=\sum_i\alpha_i y_i x_i,\qquad
\sum_i\alpha_i y_i=0,\qquad
\alpha_i = C-\mu_i .
\end{equation}
\paragraph{Dual.} En réinjectant,
\begin{equation}
\max_{\alpha}\ \sum_i\alpha_i-\tfrac12\sum_{i,j}\alpha_i\alpha_j\,y_iy_j\,
\inner{x_i}{x_j}
\quad\text{s.c.}\quad 0\le\alpha_i\le C,\ \sum_i\alpha_i y_i=0 .
\end{equation}
Problème quadratique concave où les données n'interviennent \emph{que} par
produits scalaires.
\paragraph{KKT et vecteurs supports.} La complémentarité
$\alpha_i\big[y_i(w\T x_i+b)-1+\xi_i\big]=0$ impose $\alpha_i\neq0$ uniquement pour
les points sur ou dans la marge — les \textbf{vecteurs supports} : la solution
n'en dépend que d'eux (96\,\% des points ici, donnée difficile).
\paragraph{Kernel trick.} Tout passant par $\inner{x_i}{x_j}$, on substitue un
noyau $k(x_i,x_j)=\inner{\varphi(x_i)}{\varphi(x_j)}$. Pour un noyau symétrique
défini positif, une telle application $\varphi$ existe (théorème de Mercer), et la
solution s'écrit $w=\sum_i\alpha_i y_i\varphi(x_i)$, d'où le classifieur
$f(x)=\sum_i\alpha_i y_i\,k(x_i,x)+b$. Le RBF
$k(x,x')=\exp(-\gamma\norm{x-x'}^2)$ correspond à un $\varphi$ de dimension
infinie.

\subsection{Passage à l'échelle : approximation de Nyström}\label{sec:nystrom}
La matrice de Gram exacte $K\in\R^{n\times n}$ est inabordable pour $n=50\,000$.
On l'approche par une matrice de \textbf{rang faible} $m\ll n$. Soit $m$ points de
référence (\emph{landmarks}) et les blocs
\[
C\in\R^{n\times m}\ (\text{noyau tous}\times\text{landmarks}),\qquad
W\in\R^{m\times m}\ (\text{landmarks}\times\text{landmarks}),
\]
\begin{equation}
\boxed{\ \tilde K = C\,W^{+}C\T\ }\qquad(\text{rang}\le m),
\end{equation}
$W^{+}$ étant le pseudo-inverse. \textbf{Carte de features explicite :}
diagonalisons $W=U\Lambda U\T$ et posons $Z=C\,U\,\Lambda^{-1/2}\in\R^{n\times m}$.
Alors
\begin{equation}
ZZ\T = C\,U\Lambda^{-1}U\T C\T = C\,W^{+}C\T = \tilde K ,
\end{equation}
si bien que $\inner{z_i}{z_j}\approx k(x_i,x_j)$ : on a \emph{rendu le noyau
explicite}. Le SVM à noyau devient un SVM \emph{linéaire} dans $\R^{m}$ (coût
$O(nm)$), applicable aux $50\,000$ images sans jamais former $K$.

\paragraph{Pourquoi cela fonctionne.} Si le spectre de $K$ décroît vite (cas du
RBF), $K$ est presque de rang faible et $\tilde K$ en est une bonne
reconstruction, comme pour Eckart–Young en PCA (\S\ref{sec:pca}). L'erreur
relative mesurée $\norm{K-\tilde K}_F/\norm{K}_F$ vaut $0.167,\ 0.093,\ 0.059$
pour $m=50,150,400$ : elle chute vite, ce qui confirme cette décroissance du
spectre.

\subsection{Résultats et choix}
\begin{center}
\begin{tabular}{lcc}
\toprule
Modèle & données & acc.\ test \\
\midrule
$k$-NN ($k{=}10$)        & 10k/2k & 0.281 \\
Softmax                  & 10k/2k & 0.292 \\
SVM linéaire (hinge)     & 10k/2k & 0.308 \\
SVM RBF \emph{exact}     & 10k/2k & 0.479 \\
Nyström RBF $m{=}500$    & \textbf{50k}/10k & 0.481 \\
Nyström RBF $m{=}1000$   & \textbf{50k}/10k & \textbf{0.504} \\
\bottomrule
\end{tabular}
\end{center}
\noindent\textbf{Choix.} Le RBF est retenu (non-linéarité nécessaire : le linéaire
plafonne à $0.31$). Nyström sur $50\,000$ atteint la précision du RBF exact sur
$10\,000$ \emph{pour un coût comparable} : à budget égal, il vaut mieux plus de
données via une approximation contrôlée que le noyau exact sur peu de données. En
pixels bruts on plafonne toutefois vers $0.50$ : la limite est la
\emph{représentation}, d'où l'exercice~7.

\section{Exercice 7 — Représentations non supervisées}
On apprend une application $\varphi$ \emph{sans étiquettes}, puis on teste les
latents par une \textbf{sonde linéaire} (classifieur linéaire sur représentations
gelées, $10\,000$ exemples d'entraînement, test complet). Une bonne représentation
rend la tâche \emph{linéairement} séparable.

\paragraph{Choix du scoring.} L'étape non supervisée (PCA, $K$-means) n'utilise pas
les étiquettes ; on ne peut donc pas la noter directement par une précision. On
choisit comme score la \textbf{précision de la sonde linéaire} : elle mesure
exactement ce qui nous intéresse, à savoir si la représentation rend les classes
\emph{linéairement} séparables, et reste comparable d'une représentation à l'autre
(même classifieur, même taille d'entraînement). Une mesure non supervisée interne
comme l'inertie du $K$-means renseignerait sur la compacité des amas mais pas sur
leur utilité pour la classification, qui est notre but.

\subsection{PCA : dérivation (Eckart–Young)}\label{sec:pca}
Sur données centrées, covariance $\Sigma=\tfrac1n X\T X=U\Lambda U\T$,
$\lambda_1\ge\cdots\ge\lambda_d\ge0$. Deux lectures équivalentes :
\emph{variance maximale} ($\arg\max_{\norm u=1} u\T\Sigma u = u_1$, quotient de
Rayleigh) et \emph{reconstruction minimale}
($\min \sum_i\norm{x_i-Px_i}^2$ sur les projecteurs de rang $k$), donnant le même
$U_k=[u_1,\dots,u_k]$.
\paragraph{Théorème d'Eckart–Young.} $X_k=XU_kU_k\T$ est la meilleure
approximation de rang $k$ de $X$ (normes de Frobenius et spectrale), d'erreur
$\sum_{j>k}\lambda_j$. C'est ce qui justifie de garder les premières composantes
comme représentation. Mesuré : variance cumulée $0.86$ à $k{=}64$, $0.96$ à
$k{=}256$.

\subsection{Dictionnaire $K$-means (Coates \& Ng)}
$K$-means apprend $K$ centroïdes minimisant $\sum_i\min_k\norm{p_i-c_k}^2$
(algorithme de Lloyd : assignation $\leftrightarrow$ moyenne, décroissance
monotone, convergence vers un minimum local) : une quantification vectorielle
servant de \emph{dictionnaire}. La métrique est la \textbf{distance euclidienne},
cohérente avec l'objectif quadratique de $K$-means et appliquée \emph{après}
blanchiment ZCA, qui décorrèle les pixels pour que cette distance ne soit pas
faussée par leurs corrélations. Pipeline :
\begin{enumerate}\itemsep1pt
  \item patches $6\!\times\!6\!\times\!3$ aléatoires, normalisés en contraste
  (centrage + variance unité) ;
  \item \textbf{blanchiment ZCA} $p\mapsto Wp$, $W=U(\Lambda+\varepsilon I)^{-1/2}U\T$
  (covariance des patches) : décorrélation — encore le théorème spectral ;
  \item dictionnaire : $K$-means sur patches blanchis ;
  \item encodage \emph{triangle}
  $f_k(p)=\max\!\big(0,\bar d(p)-\norm{p-c_k}\big)$ (code parcimonieux, non négatif),
  $\bar d$ distance moyenne aux centroïdes ;
  \item \textbf{pooling} par quadrants $2\!\times\!2 \Rightarrow 4K$ features,
  passées à un SVM linéaire.
\end{enumerate}
Intuition : blanchiment + $K$-means fabriquent des détecteurs de bords/couleurs ;
le pooling apporte l'invariance ; la tâche devient linéairement séparable.

\subsection{Résultats et choix}
\begin{center}
\begin{tabular}{lcc}
\toprule
Représentation & dim & acc.\ test (sonde lin.) \\
\midrule
Pixels bruts                       & 3072 & 0.323 \\
PCA                                & 256  & 0.380 \\
\textbf{Dictionnaire $K$-means}    & 1024 & \textbf{0.659} \\
\bottomrule
\end{tabular}
\end{center}
\noindent\textbf{Choix.} À classifieur \emph{linéaire fixé}, passer des pixels au
dictionnaire $K$-means double quasiment la précision ($0.32\!\to\!0.66$), et sans
étiquette : la difficulté de CIFAR est surtout dans la représentation, pas dans le
classifieur. Augmenter $K$ (jusqu'à $1600$) pousse vers $\sim0.77$ (Coates \& Ng) ;
on garde $K{=}256$ pour un temps de calcul modéré.

\section{Conclusion}
On retrouve deux outils principaux. D'une part l'optimisation convexe : le SVM se
résout via son dual (KKT, vecteurs supports), le kernel trick déplace la
non-linéarité, et Nyström permet le passage à l'échelle. D'autre part la
diagonalisation des matrices symétriques (théorème spectral), qui sert pour la PCA,
le blanchiment ZCA et le bloc $W$ de Nyström. Le gain décisif vient finalement de
l'apprentissage non supervisé de représentations ($K$-means, PCA) plutôt que du
choix du classifieur.

\end{document}
